\section{Quinta iteración de la Fase 5: Validación pasa al nuevo módulo de cartuchos}
\label{sec:F5I5}

La verificación anterior al guardado del GWI — la coherencia
de las listas que el diseñador deja en memoria — no es lógica del
editor: es el \textbf{contrato del formato}, lo que un cartucho
debe cumplir para poder escribirse, igual que \texttt{validarCartucho}
define lo que puede leerse. Esta iteración es un
\textbf{traspaso puro}, como el de la anterior: mudar desde 
\texttt{BaseDescomponer.js} a \texttt{BaseCartucho.js}, sin cambiar
su comportamiento, pero aprovechar para darle la forma que sus dos
futuros puntos de llamada necesitan: una \textbf{función independiente}
que devuelve su resultado en vez de abrir el guardado.

\subsection{Qué se traspasa}
\label{sec:F5I5Que}

La verificación completa del diseñador: el recorrido de las
listas que exige que cada cartucho que va a salir tenga sentido
por sí solo. Comprueba que los enlaces apunten a segmentos
válidos y con la dirección correcta —calculando la inversa si
falta, como las direcciones—, que las colisiones estén declaradas
con sus contadores de desenlace, que los secundarios que se describen en la 
Sección~\ref{sec:SegmentosPresentanJuntos} (treceava iteración de la fase 3 para los
segmentos que se presentan juntos), no formen red de enlaces ni queden huérfanos,
que los contadores de dígitos estén completos, que exista un
jugador inicial, y que el fin de partida y la Secuencia estén
declarados con sentido. Su resultado ya era de dos clases, y
sigue siéndolo: los \textbf{errores}, que impiden guardar, y los
\textbf{avisos}, que solo informan, este reparto no cambia.

Todo se traslada tal cual funciona actualmente, con la adaptación necesaria para
que la verificación reciba las listas por parámetro: el módulo
no conoce al editor, así que no puede leer su memoria interna, usando 
el mismo criterio que \texttt{validarCartucho} ya aplica al
recibir el estilo de la máquina.

\subsection{La función independiente}
\label{sec:F5I5Funcion}

Lo que hoy es el preámbulo de «Verificar y Guardar» pasa a ser
una función propia del módulo con esta firma:

		
\begin{description}
	\item \texttt{verificar(listas)} \\
	Comprueba la coherencia del cartucho antes de escribirlo: recorre las listas, calcula lo que falte y devuelve los errores que impiden guardar y los avisos que solo informan.
\end{description}


Los procesos de guardado la llaman \textbf{siempre primero}, y la
severidad la decide el punto de llamada: hoy solo «Verificar y
Guardar» existe, y es bloqueante con los errores; cuando la
siguiente iteración añada «Generar GWI+GWC», llamará a esta
misma función antes que nada, sin duplicar el código y sin necesidad de validar dos
veces lo mismo. Junto a ella, ya en su sitio, queda también
\texttt{validarTopes}, que llegará con el GWC: cierto o falso, y
el aviso o el bloqueo según quién la llame.

\subsection{Los cambios en los módulos}
\label{sec:F5I5Cambios}

La mudanza es de traspaso, no de reescritura:

\begin{itemize}
	\item \textbf{BaseCartucho} gana la sección de validación de
	escritura, junto a la validación de lectura que ya tenía: todo
	lo que decide si un cartucho puede entrar o salir vive en el
	mismo módulo.
	
	\item \textbf{BaseDescomponer} pierde la verificación y queda
	con su llamada: el botón «Verificar y Guardar» llama a
	\texttt{BaseCartucho.verificar}, y solo si no hay errores
	continúa con el montaje y el guardado de siempre. Se quedan en
	él lo que es del editor: el check del cartucho en edición, el
	montaje del GWI y el repintado.
	
	\item \textbf{index.html} no cambia: el módulo ya se carga
	desde el traspaso.
\end{itemize}

Ningún otro módulo cambia: el motor, el puente y el generador
gráfico quedan exactamente como están.

\subsection{Pruebas del traspaso}
\label{sec:F5I5Pruebas}

El traspaso es puro, así que la prueba es una sola:
\textbf{el comportamiento no cambia}. En simulación y en
navegador real:

\begin{itemize}
	\item \textbf{Cartucho válido:} el GWI real del Parachute pasa
	la verificación con sus avisos de siempre y guarda igual.
	
	\item \textbf{Errores:} los mismos defectos de siempre —enlaces
	colgantes, colisiones sin contador, secundarios con red— dan
	los mismos mensajes de error, ahora saliendo del módulo.
	
	\item \textbf{Avisos:} los avisos no bloquean: el cartucho se
	guarda igual, con su cuenta de avisos en el mensaje de
	confirmación.
	
	\item \textbf{Cálculos:} las inversas y los lados que la
	verificación completa sobre la marcha siguen calculándose
	igual, y los huérfanos se retiran igual.
	
	\item \textbf{El botón:} «Verificar y Guardar» se comporta
	idéntico: primero la verificación, y sin errores el guardado
	en ZIP con el mismo contenido de siempre.
\end{itemize}

